%%%%%%%%%%%%%%%%%%%%%%%%%%%%%%%%%%%%%%%%%%%%%%%%%%%%%%%%%%%%%%%%%%%%%%%%%%%%%%%
%% Capítulo 5 --- Conclusões
%%
%% Par do capitulo01.tex: a Introdução estabelece a dívida, este capítulo a paga.
%% A afirmação central, a ordem de precedência da cadeia e a declaração de
%% procedência em três níveis são as mesmas do Capítulo 1, com a mesma força.
%%
%% Classe esperada: \documentclass[tcc,alf]{tcc-poli}
%% Citações: abntex2cite (sistema alfabético).
%% Pacotes usados: siunitx, hyperref (ambos já no principal.tex).
%%
%% MARCAÇÃO DE TRECHO EM ANDAMENTO --------------------------------------------
%% \emandamento{...} destaca em vermelho os trechos cuja validade muda quando o
%% protótipo for construído e ensaiado em ambiente real. Ver nota no capitulo01.
%%%%%%%%%%%%%%%%%%%%%%%%%%%%%%%%%%%%%%%%%%%%%%%%%%%%%%%%%%%%%%%%%%%%%%%%%%%%%%%

\makeatletter
\@ifundefined{emandamento}{%
  \@ifundefined{textcolor}%
    {\newcommand{\emandamento}[1]{#1}}%
    {\newcommand{\emandamento}[1]{\textcolor{red}{#1}}}%
}{}
\makeatother

\chapter{Conclusões}
\label{cap:conclusoes}

Este trabalho perguntou o que a cadeia de medição especificada para o Kaelix permite
medir, e não se um dispositivo de baixo custo pode ser construído. A resposta é que a
cadeia sustenta a detecção de desbalanceamento e de desvios de condição de operação sob
taxa de falso alarme fixada, e não sustenta diagnóstico de defeito de rolamento. O limite
não decorre de escolha de processamento e não é removível por ajuste de \emph{software}:
ele está na aquisição, e a ordem de precedência entre os elos da cadeia determina que
nenhuma decisão posterior recupere a informação que o sensor não coletou.

A contribuição verificável deste trabalho está no segundo elo. As decisões de limiar, de
métrica e de protocolo de validação alteraram os resultados por margens maiores que o
efeito físico sob investigação, e é sobre elas que os experimentos produziram números
novos. Essa contribuição depende de que o primeiro elo esteja declarado: um desempenho de
detecção reportado sem a banda e o piso de ruído que o produziram não é verificável, e a
prática de omiti-los é o que a revisão encontrou de mais frágil na literatura de vibração
embarcada.

\section{O que a aquisição permite medir}
\label{sec:conc-aquisicao}

A taxa máxima de amostragem de \SI{1}{\kilo\hertz} do acelerômetro situa a frequência de
Nyquist em \SI{500}{\hertz}. A frequência característica de defeito de pista externa, de
\SI{104,6}{\hertz} na geometria considerada, cai dentro dessa banda; a ressonância
estrutural de \SI{3,5}{\kilo\hertz} que carrega a informação diagnóstica não. Essa
separação é o resultado central do primeiro elo: a taxa de repetição do defeito estar
dentro da banda não significa que a evidência do defeito esteja.

Sobre o sinal simulado, o contraste entre as condições com e sem defeito foi de
\num{413} vezes por análise de envelope na banda de ressonância, amostrada a
\SI{20}{\kilo\hertz}, contra \num{20} vezes por espectro direto a \SI{1}{\kilo\hertz}.
Cabe a qualificação que a revisão já impõe a esse par: entre as duas medidas mudam duas
variáveis, a banda e o método de demodulação, e dois pontos não constituem uma função da
banda. O par sustenta a afirmação de que a cadeia truncada perde contraste; não sustenta a
atribuição quantitativa desse ganho à banda isoladamente. Separá-los exige varredura com
método fixo, e é trabalho futuro (Seção~\ref{sec:conc-futuro}).

Um segundo efeito é de natureza distinta e merece registro por ser, ele próprio, um
achado. Sem filtro passa-baixas configurado abaixo de $f_s/2$, a ressonância de
\SI{3,5}{\kilo\hertz} não desaparece: reaparece rebatida em \SI{500}{\hertz}, uma
frequência onde nada existe fisicamente. Descritores estatísticos medidos sobre esse sinal
podem separar as classes por um mecanismo que nada tem a ver com a física do defeito. Uma
medição ingênua pode, portanto, funcionar pelos motivos errados e passar em validação.

O invólucro integra a cadeia de medição, e não apenas a protege. A rigidez do caminho entre
a superfície do motor e o sensor situa a frequência de montagem do conjunto em
\SI{557}{\hertz}, dentro da banda de medição, com transmissibilidade de aproximadamente
cinco vezes em \SI{500}{\hertz} e erro de amplitude superior a \SI{10}{\percent} já em
\SI{168}{\hertz}. A prática consolidada em acelerometria exigiria frequência de montagem de
ao menos \SI{3}{\kilo\hertz} para a banda normativa. A decomposição da flexibilidade
localiza o problema na base e no corpo, responsáveis por \SI{51,1}{\percent} e
\SI{48,5}{\percent}, e não no ressalto, com \SI{0,4}{\percent}; a mesma geometria em
alumínio alcançaria \SI{3273}{\hertz}. Duas hipóteses concorrentes foram testadas e
descartadas: a curvatura da carcaça e os modos próprios dos painéis, estes situados acima
de \SI{1300}{\hertz}.

\section{O que a estimação decide}
\label{sec:conc-estimacao}

Três decisões do segundo elo alteraram os resultados por margens maiores que o efeito
medido.

A primeira é o método de integração. A norma define as zonas de severidade em velocidade
eficaz, e o acelerômetro entrega aceleração. Contra um sinal cuja integral é conhecida
analiticamente, a integração no domínio do tempo errou em $+\SI{73}{\percent}$ sem filtro
passa-alta e em $+\SI{156}{\percent}$ com ele, enquanto a integração no domínio da
frequência, restrita à máscara normativa, recuperou o valor exato. O filtro passa-alta não
corrige a integração no tempo e chega a piorá-la. Sob viés de
\SI{0,02}{\meter\per\second\squared}, valor contido na especificação do sensor, a
integração no tempo produziu \SI{55,08}{\milli\meter\per\second} onde o correto é
\num{6,75}: um sinal da zona A da norma seria classificado além da zona D. A integração na
frequência permanece exata sob o mesmo viés, porque o termo contínuo cai fora da máscara
imposta pela própria norma.

A segunda é a política de limiar. Treinado apenas com dados de operação normal, o detector
com o valor padrão do parâmetro de contaminação classificou \SI{42,1}{\percent} das
janelas normais como anomalia; com o parâmetro fixado em \num{0,01}, o falso alarme caiu
para \SI{1,4}{\percent}. A taxa de detecção permaneceu em \SI{100}{\percent} e a área sob
a curva ROC em \num{1,000} nos quatro valores testados. O parâmetro desloca o corte e não
altera a qualidade do detector, o que confirma que o desempenho reportado sem a política
de limiar declarada não é interpretável.

A terceira é a métrica. Sobre defeito incipiente, com descritores e modelo idênticos, a
área sob a curva completa caiu \SI{3,7}{\percent} entre a cadeia de referência e a do
dispositivo, enquanto a área parcial restrita à região de falso alarme baixo caiu
\SI{16,2}{\percent} a \SI{10}{\percent} de falso alarme e \SI{20,3}{\percent} a
\SI{5}{\percent}. A perda na região em que um dispositivo de campo precisa operar é quatro
a cinco vezes maior do que a área global indica.

A essas três soma-se o protocolo de validação. Janelas consecutivas de um mesmo ensaio
compartilham a assinatura da montagem, e sua divisão aleatória entre treino e teste permite
que o modelo identifique o ensaio em vez da falha. O mesmo modelo, sobre os mesmos dados e
os mesmos descritores, atingiu \SI{99,0}{\percent} de acurácia sob divisão aleatória por
janela e $\num{63,2} \pm \num{19,2}\,\si{\percent}$ sob divisão por ensaio, uma inflação de
aproximadamente \num{36} pontos percentuais. Três das cinco partições ficaram próximas de
\num{0,50}, o nível do acerto ao acaso, informação que a média única do protocolo
incorreto suprime.

\section{O que a implantação impõe}
\label{sec:conc-implantacao}

A hipótese inicial atribuía ao material do invólucro o papel de elo térmico mais frágil. A
análise inverteu essa leitura. A resistência térmica do caminho condutivo em ASA é de
\SI{88,8}{\kelvin\per\watt}, contra \SI{0,074}{\kelvin\per\watt} em alumínio, e com a
carcaça a \SI{80}{\celsius} a temperatura interna estabiliza em \SI{44,8}{\celsius}. O
autoaquecimento é irrelevante: a potência média de \SI{2,0}{\milli\watt} eleva a
temperatura interna em \SI{0,01}{\celsius}, de modo que todo o calor provém do exterior.

O componente que atinge primeiro seu limite é a bateria, pela restrição de carga em
\SI{45}{\celsius}, temperatura interna alcançada com carcaça em torno de
\SI{81}{\celsius}. O limite de operação do dispositivo é imposto pela célula, e não pelo
invólucro. A consequência de projeto é que substituir o material do invólucro por alumínio,
que resolveria a limitação mecânica da Seção~\ref{sec:conc-aquisicao}, agravaria a
limitação térmica, porque elevaria a temperatura interna a \SI{79,0}{\celsius} sob a mesma
carcaça. Os dois elos impõem requisitos opostos sobre a mesma peça.

\section{Alcance dos resultados}
\label{sec:conc-alcance}

As conclusões acima têm alcances distintos, porque a procedência dos resultados não é
uniforme, e enunciá-las sem essa distinção seria atribuir a todas a força da mais forte.

Os resultados de detecção em modos de falha provêm de dado real medido em bancada
instrumentada, submetido a decimação controlada em \emph{software} para emular a banda do
dispositivo: a \SI{2}{\percent} de falso alarme, \SI{75,0}{\percent} de detecção em
desbalanceamento, entre \num{15} e \SI{21}{\percent} em defeito de rolamento puro e entre
\num{5} e \SI{7}{\percent} em desalinhamento. O que é modelado nesses resultados é a
aquisição, não o sinal.

Os resultados de banda, de integração, de limiar e de protocolo provêm de sinais
sintetizados a partir de um modelo de dois mecanismos, e valem para a cadeia de
processamento, não para o desempenho em campo. Seus valores absolutos dependem dos
parâmetros de simulação --- amplitude do defeito, frequência de ressonância e nível de
ruído --- e devem ser lidos como ordens de grandeza. As relações entre métodos, que são o
objeto deste trabalho, não dependem desses parâmetros.

Os resultados mecânicos e térmicos provêm de modelagem analítica sobre as especificações
de projeto, tomadas como dadas. A análise mecânica emprega fórmulas fechadas com condições
de contorno idealizadas, e a rigidez real depende de pré-carga de parafusos, planicidade
das faces impressas e orientação das camadas. A análise térmica assume regime permanente e
não modela resistência de contato, convecção forçada pelo ventilador do motor nem radiação
de superfícies próximas.

\emandamento{A cadeia física do próprio dispositivo não foi verificada. Nenhum ensaio
físico foi conduzido e nenhuma medição de campo alimenta os modelos; a inicialização do
acelerômetro permanece não implementada no firmware e o filtro passa-baixas interno do
sensor é pendência declarada em código. A consequência é uma restrição específica, e não
uma ressalva genérica: os resultados sustentam conclusões sobre o que a cadeia de
aquisição especificada permite medir, e não sobre o desempenho do dispositivo montado em
operação industrial.} Em particular, a bancada sobre a qual os dados reais foram obtidos está abaixo
do escopo de potência e altura de eixo da norma, o que a torna adequada como fonte de
sinal para estudar a cadeia de medição e inadequada como caso de aplicação normativa.

\emandamento{Três dados externos permanecem pendentes e limitam o alcance das conclusões:
a temperatura real de carcaça dos equipamentos-alvo, a classe normativa dos motores a
monitorar e a validação dos carregadores dos conjuntos de dados públicos utilizados.}

\section{Trabalho futuro}
\label{sec:conc-futuro}

\emandamento{A construção do protótipo físico e seu ensaio em ambiente real constituem a
etapa seguinte declarada deste trabalho, e não uma possibilidade remota. Os resultados
apresentados aqui não substituem essa etapa: eles a preparam, ao estabelecer o que ela
precisa verificar e quais medidas tornam essa verificação interpretável. Um ensaio de
campo conduzido sem banda declarada, sem política de limiar explícita e sob divisão
aleatória de janelas reproduziria, em dado real, os mesmos artefatos que este trabalho
mediu em sinal controlado.}

Cada item abaixo é enunciado junto da afirmação que ele permitiria sustentar, e não como
lista de desejos.

\begin{itemize}
    \item \textbf{Varredura de taxa de amostragem com método fixo}, sobre a mesma gravação
          decimada em vários $f_s$, com decomposição entre efeito de banda e efeito de
          perda do envelope. É o que converte o par \num{413}$\times$/\num{20}$\times$ em
          função da banda, e o que falta para atribuir a perda de contraste à banda
          isoladamente.

    \item \textbf{Medição dos mesmos descritores com e sem filtro passa-baixas
          configurado}, no mesmo sinal. É o que estabelece experimentalmente o artefato de
          rebatimento sobre curtose e fator de crista, efeito para o qual não se localizou
          descrição na literatura revisada por pares.

    \item \textbf{Varredura do parâmetro de contaminação contra falso alarme medido}, em
          vibração industrial de banda limitada. É o que permite afirmar que a calibração
          por risco alvo supera a política padrão, em vez de apenas exibir que a política
          padrão falha.

    \item \textbf{Taxa de falso alarme de detector treinado no próprio microcontrolador.}
          É o modo de operação exigido de uma máquina sem histórico rotulado, e nenhum
          valor publicado foi localizado para ele.

    \item \textbf{Descritores de razão harmônica para desalinhamento.} O desempenho em
          desalinhamento, entre \num{5} e \SI{7}{\percent}, é o pior dos três modos
          medidos e não é explicável por banda: o modo se manifesta em harmônicas da
          rotação, muito abaixo de \SI{500}{\hertz}. A causa provável está nos descritores
          adotados, que capturam energia e impulsividade e não estrutura harmônica e de
          fase. É a pista mais acionável do conjunto, e não custa banda, sensor nem
          dinheiro.

    \item \textbf{Medição por excitação de impacto da frequência natural do conjunto
          montado}, incluindo a ressonância da fixação magnética. É o que substitui a
          estimativa analítica de \SI{557}{\hertz}, que depende de um amortecimento
          arbitrado, por valor medido.

    \item \emandamento{\textbf{Construção do protótipo e ensaio em ambiente real} ---
          etapa seguinte já prevista. É a única via que converte as conclusões sobre a
          cadeia especificada em conclusões sobre o dispositivo em operação. Três
          verificações a precedem no próprio dispositivo: implementar a inicialização do
          acelerômetro, configurar o filtro passa-baixas interno abaixo de $f_s/2$ e medir
          o piso de ruído efetivo, em vez de adotar o valor de folha de dados. Sem as duas
          primeiras, a medição de campo mede aliasing; sem a terceira, não há como separar
          defeito não detectado de defeito abaixo do piso.}
\end{itemize}

A ordem entre esses itens não é arbitrária. Os dois primeiros pertencem ao primeiro elo da
cadeia e condicionam a interpretação de todo o resto; os três seguintes pertencem ao
segundo elo, onde este trabalho localizou sua contribuição; o penúltimo pertence ao
terceiro. O último não é um item entre outros: é a continuação do trabalho, e o que este
documento oferece a ela é o conjunto de medidas que tornam seu resultado interpretável.

\emandamento{Nesse sentido, as conclusões aqui apresentadas são de uma etapa, e não de um
ciclo fechado. Elas delimitam o que a especificação permite esperar do dispositivo antes
que ele exista fisicamente, o que é precisamente o que permite projetar o ensaio de campo
de modo a testar afirmações, e não a confirmá-las.}
